\documentclass[10pt,a4paper]{amsart}
\usepackage{amsmath,amssymb,amsthm,graphicx}
% Retained minimal preamble for the Stacks Project chapter more-algebra.tex (More on Algebra).
% The chapter file pinned-original-source.tex carries only \input{preamble} on its first line, so no
% preamble text lives inside the chapter itself.  The Stacks Project uses one single shared preamble for
% every chapter; the snapshot stages it identically as sources/stacks-chow/preamble.tex and
% sources/stacks-etale-cohomology/preamble.tex.  The block below is copied mechanically, line for line,
% from sources/stacks-chow/preamble.tex lines 14-18 (the Xy-pic setup) and 154-217 (the native theorem
% declarations together with the mathematical macros).  Nothing was added, removed or redefined; the
% wrapper class and the amsmath/amssymb/amsthm/graphicx base packages are supplied by the assembler.
% Chapter: more-algebra.tex -- More on Algebra.
% For commutative diagrams we use Xy-pic
\usepackage[all]{xy}

% We use 2cell for 2-commutative diagrams.
\xyoption{2cell}

% Theorem environments.
%
\theoremstyle{plain}
\newtheorem{theorem}[subsection]{Theorem}
\newtheorem{proposition}[subsection]{Proposition}
\newtheorem{lemma}[subsection]{Lemma}

\theoremstyle{definition}
\newtheorem{definition}[subsection]{Definition}
\newtheorem{example}[subsection]{Example}
\newtheorem{exercise}[subsection]{Exercise}
\newtheorem{situation}[subsection]{Situation}

\theoremstyle{remark}
\newtheorem{remark}[subsection]{Remark}
\newtheorem{remarks}[subsection]{Remarks}

\numberwithin{equation}{subsection}

% Macros
%
\def\lim{\mathop{\mathrm{lim}}\nolimits}
\def\colim{\mathop{\mathrm{colim}}\nolimits}
\def\Spec{\mathop{\mathrm{Spec}}}
\def\Hom{\mathop{\mathrm{Hom}}\nolimits}
\def\Ext{\mathop{\mathrm{Ext}}\nolimits}
\def\SheafHom{\mathop{\mathcal{H}\!\mathit{om}}\nolimits}
\def\SheafExt{\mathop{\mathcal{E}\!\mathit{xt}}\nolimits}
\def\Sch{\mathit{Sch}}
\def\Mor{\mathop{\mathrm{Mor}}\nolimits}
\def\Ob{\mathop{\mathrm{Ob}}\nolimits}
\def\Sh{\mathop{\mathit{Sh}}\nolimits}
\def\NL{\mathop{N\!L}\nolimits}
\def\CH{\mathop{\mathrm{CH}}\nolimits}
\def\proetale{{pro\text{-}\acute{e}tale}}
\def\etale{{\acute{e}tale}}
\def\QCoh{\mathit{QCoh}}
\def\Ker{\mathop{\mathrm{Ker}}}
\def\Im{\mathop{\mathrm{Im}}}
\def\Coker{\mathop{\mathrm{Coker}}}
\def\Coim{\mathop{\mathrm{Coim}}}

% Boxtimes
%
\DeclareMathSymbol{\boxtimes}{\mathbin}{AMSa}{"02}

%
% Macros for moduli stacks/spaces
%
\def\QCohstack{\mathcal{QC}\!\mathit{oh}}
\def\Cohstack{\mathcal{C}\!\mathit{oh}}
\def\Spacesstack{\mathcal{S}\!\mathit{paces}}
\def\Quotfunctor{\mathrm{Quot}}
\def\Hilbfunctor{\mathrm{Hilb}}
\def\Curvesstack{\mathcal{C}\!\mathit{urves}}
\def\Polarizedstack{\mathcal{P}\!\mathit{olarized}}
\def\Complexesstack{\mathcal{C}\!\mathit{omplexes}}
% \Pic is the operator that assigns to X its picard group, usage \Pic(X)
% \Picardstack_{X/B} denotes the Picard stack of X over B
% \Picardfunctor_{X/B} denotes the Picard functor of X over B
\def\Pic{\mathop{\mathrm{Pic}}\nolimits}
\def\Picardstack{\mathcal{P}\!\mathit{ic}}
\def\Picardfunctor{\mathrm{Pic}}
\def\Deformationcategory{\mathcal{D}\!\mathit{ef}}

% --- appended 2026-10-12 (Chief): the project own definitions of the reference/slogan/history
% display environments, copied mechanically from the shared Stacks preamble lines 8-12
% (\usepackage{verbatim} plus the three \newenvironment lines, which the chapter preamble keeps
% verbatim from the project preamble). Nothing added, removed or redefined: these environments
% carry no mathematical content and the project itself typesets them as comments. Without these
% three lines a counted statement whose own text uses them fails with Environment slogan undefined.
\usepackage{verbatim}
\newenvironment{reference}{\comment}{\endcomment}
\newenvironment{slogan}{\comment}{\endcomment}
\newenvironment{history}{\comment}{\endcomment}
\usepackage{cleveref}
\crefname{originallemma}{original lemma}{original lemmas}
\Crefname{originallemma}{Original lemma}{Original lemmas}
\makeatletter
\newlabel{more-algebra-lemma-flat-graded-finite-type-finite-presentation-module}{{053C}{source}{}{sourcepointer-1}{}}
\newlabel{more-algebra-lemma-flat-graded-finite-type-finite-presentation-module@cref}{{[originallemma][0][]053C}{[1][0][]source}{}{}{}}
\makeatother
\crefname{originallemma}{original lemma}{original lemmas}
\Crefname{originallemma}{Original lemma}{Original lemmas}
\makeatletter
\newlabel{algebra-lemma-finite-flat-local}{{00NZ}{source}{}{sourcepointer-2}{}}
\newlabel{algebra-lemma-finite-flat-local@cref}{{[originallemma][0][]00NZ}{[1][0][]source}{}{}{}}
\makeatother
\crefname{originallemma}{original lemma}{original lemmas}
\Crefname{originallemma}{Original lemma}{Original lemmas}
\makeatletter
\newlabel{algebra-lemma-flat-ses}{{00HM}{source}{}{sourcepointer-3}{}}
\newlabel{algebra-lemma-flat-ses@cref}{{[originallemma][0][]00HM}{[1][0][]source}{}{}{}}
\makeatother
\crefname{originallemma}{original lemma}{original lemmas}
\Crefname{originallemma}{Original lemma}{Original lemmas}
\makeatletter
\newlabel{algebra-lemma-NAK}{{00DV}{source}{}{sourcepointer-4}{}}
\newlabel{algebra-lemma-NAK@cref}{{[originallemma][0][]00DV}{[1][0][]source}{}{}{}}
\makeatother
\title{A finite graded module over a graded polynomial algebra that is flat over a local ring is finitely presented}
\author{The Stacks Project Authors}
\date{stacks-project:more-algebra.tex:local-snapshot-20261012; GNU FDL 1.2 or later}
\begin{document}
\maketitle
\setcounter{subsection}{1}
\noindent\textit{Editorial note.} All delivered mathematical text is the author's own, copied from the pinned Stacks Project chapter file more-algebra.tex; no statement or proof step was rewritten, reordered or removed.\par\medskip
\noindent\textit{Editorial note.} The chapter file pinned-original-source.tex carries only \textbackslash{}input\{preamble\} on its first line, so no preamble text lives inside the chapter. The retained file sources/stacks-more-algebra/retained-preamble.tex is a mechanical, line-for-line copy of the project's single shared preamble as staged in the snapshot (Xy-pic setup plus theorem declarations and mathematical macros). Nothing was added, removed or redefined, and the excerpt serves only as preamble, never as a span coordinate.\par\medskip
\noindent\textit{Editorial note.} Both counted spans start at the \textbackslash{}begin line of their environment, so the statement and the author's proof are delivered whole; surrounding prose of the section is deliberately outside the counted spans and is carried as standard prerequisites instead.\par\medskip
\noindent\textit{Editorial note.} The citation fields (source\_id, source\_version, authors, source\_url, license.name, source\_title) are identical to those of every other card this lane has produced from this chapter, so the batch selector treats them as one source.\par\medskip
\noindent\textit{Allowed standard prerequisites (editorial).}\begin{itemize}
\item A finite flat module over a local ring is free. Algebra, Lemma \textbackslash{}ref\{algebra-lemma-finite-flat-local\}, cited at proof lines 5884 and 5899. Makes each graded piece of the module, and then each graded piece of the kernel, a finite free module.
\item Flatness passes to the kernel of a surjection of flat modules. Algebra, Lemma \textbackslash{}ref\{algebra-lemma-flat-ses\}, cited at proof lines 5895-5897. Shows that each graded piece K\_d of the kernel is flat over R.
\item Nakayama's lemma. Algebra, Lemma \textbackslash{}ref\{algebra-lemma-NAK\}, cited at proof line 5907. Turns a generating set of the kernel after tensoring with the residue field into a generating set over the ring.
\end{itemize}\medskip
\ifdefined\hypertarget\hypertarget{sourcepointer-1}{}\fi
\noindent\textit{Source reference.} Original lemma 053C: focused-seeds/AGP/sources/stacks-snapshot-20261011/tags.txt, source lines 6600--6600. The lemma's own canonical tag, quoted so that readers can resolve the label inside the delivered text; it points at the counted statement itself and carries no proof content.\par\medskip
\ifdefined\hypertarget\hypertarget{sourcepointer-2}{}\fi
\noindent\textit{Source reference.} Original lemma 00NZ: focused-seeds/AGP/sources/stacks-snapshot-20261011/tags.txt, source lines 877--877. Navigation pointer to a statement the counted proof refers to but does not contain; nothing load-bearing is delegated to it.\par\medskip
\ifdefined\hypertarget\hypertarget{sourcepointer-3}{}\fi
\noindent\textit{Source reference.} Original lemma 00HM: focused-seeds/AGP/sources/stacks-snapshot-20261011/tags.txt, source lines 648--648. Navigation pointer to a statement the counted proof refers to but does not contain; nothing load-bearing is delegated to it.\par\medskip
\ifdefined\hypertarget\hypertarget{sourcepointer-4}{}\fi
\noindent\textit{Source reference.} Original lemma 00DV: focused-seeds/AGP/sources/stacks-snapshot-20261011/tags.txt, source lines 512--512. Navigation pointer to a statement the counted proof refers to but does not contain; nothing load-bearing is delegated to it.\par\medskip
\begin{lemma}
\label{lemma-flat-graded-finite-type-finite-presentation-module}
Let $R$ be a ring.
Let $S = R[x_1, \ldots, x_n]$ be a graded polynomial algebra over $R$,
i.e., $\deg(x_i) > 0$ but not necessarily equal to $1$.
Let $M$ be a graded $S$-module.
Assume
\begin{enumerate}
\item $R$ is a local ring,
\item $M$ is a finite $S$-module, and
\item $M$ is flat over $R$.
\end{enumerate}
Then $M$ is finitely presented as an $S$-module.
\end{lemma}

\begin{proof}
Let $M = \bigoplus M_d$ be the grading on $M$.
Pick homogeneous generators $m_1, \ldots, m_r \in M$ of $M$.
Say $\deg(m_i) = d_i \in \mathbf{Z}$. This gives us a presentation
$$
0 \to K \to \bigoplus\nolimits_{i = 1, \ldots, r} S(-d_i) \to M \to 0
$$
which in each degree $d$ leads to the short exact sequence
$$
0 \to K_d \to \bigoplus\nolimits_{i = 1, \ldots, r} S_{d - d_i} \to
M_d \to 0.
$$
By assumption each $M_d$ is a finite flat $R$-module. By
Algebra, Lemma \ref{algebra-lemma-finite-flat-local}
this implies each $M_d$ is a finite free $R$-module. Hence
we see each $K_d$ is a finite $R$-module. Also each $K_d$ is flat
over $R$ by
Algebra, Lemma \ref{algebra-lemma-flat-ses}.
Hence we conclude that each $K_d$ is finite free by
Algebra, Lemma \ref{algebra-lemma-finite-flat-local}
again.

\medskip\noindent
Let $\mathfrak m$ be the maximal ideal of $R$. By the flatness of $M$
over $R$ the short exact sequences above remain short exact after tensoring
with $\kappa = \kappa(\mathfrak m)$. As the ring $S \otimes_R \kappa$ is
Noetherian we see that there exist homogeneous elements
$k_1, \ldots, k_t \in K$ such that the images $\overline{k}_j$
generate $K \otimes_R \kappa$ over $S \otimes_R \kappa$. Say $\deg(k_j) = e_j$.
Thus for any $d$ the map
$$
\bigoplus\nolimits_{j = 1, \ldots, t} S_{d - e_j}
\longrightarrow
K_d
$$
becomes surjective after tensoring with $\kappa$. By
Nakayama's lemma
(Algebra, Lemma \ref{algebra-lemma-NAK})
this implies the map is surjective over $R$. Hence $K$ is generated
by $k_1, \ldots, k_t$ over $S$ and we win.
\end{proof}

\end{document}
